% !TEX root = ./main.tex
\documentclass[a4paper]{article}
% Course language: use [english] or [german].
\usepackage[german]{ajnotes}
\title{\textbf{Betriebwirtschaftslehre — Übungseinheit C}}
\begin{document}
    \maketitle
    % Attach the given problem sheet: drop it in this folder as sheet.pdf and
    % un-comment the next line (pdfpages is already loaded by the preamble).
    % \includepdf[pages=-]{sheet.pdf}
    \section{Marketingprogramm}
    Integriertes Marketingprogramm mit konsistenter Anwendung Marketing Instrumente. Es geht um den Ps: Product und Price.

    \subsection{Produkt schaffen}

    \subsubsection{Produktdimensionen: Mahalo Wasser}
    \begin{enumerate}
        \item Kernprodukt (Inhalt)
        \begin{itemize}
            \item Aus Hawaii importiertem Wasser 
            \item Erfrischung 
            \item Erinnerung an Ferien
        \end{itemize}
        \item Reales Produkt (Form)
              \begin{itemize}
                  \item Qualität
                  \item Design
                  \item Verpackung
              \end{itemize}
        \item Erweitertes Produkt (zusätzliche Dienstleistungen)
              \begin{itemize}
                  \item Kundendienst
                  \item Kostenlose Goodies
              \end{itemize}
    \end{enumerate}

    \subsubsection{Reeles Produkt: Marke}
    \begin{itemize}
        \item Erinnerung an Vorteile des Produkts 
        \begin{itemize}
            \item Positiv: aloha erinnert an Hawaii, 
        \end{itemize}
        \item Einfache Aussprache
        \item Einzigartigkeit
        \begin{itemize}
            \item Nur geografisch begrenzt Einzigartig (z.B in der Schweiz) aber weit verbreitet als Produktmarke
        \end{itemize}
        \item Wiedererkennbarkeit
        \item Sinnvolle Übersetzung 
        \begin{itemize}
            \item Einheitliche Weltbekannte Bedeutung ohne Bedarf zur Übersetzung
        \end{itemize}
        \item Als Warenzeichen registrierbar
        \item Auf andere Produkte ausweitbar (Brand extension, Produktlinie)
        \begin{itemize}
            \item Bsp Nivea, Seife, Gesichtscreme, Deo etc.
            \item Relevant für Praxis: Marke sollte nicht auf zu vielen Produkten (zu langen Produktlinien) angewendet werden, das führt zur Verwässerung der Bedeutung der Marke.
        \end{itemize}
    \end{itemize}

    \subsubsection{Reeles Produkt: Verpackung}
    \begin{itemize}
        \item Funktionen
        \begin{itemize}
            \item Schutz
            \item Sicherheit
            \item Information
            \item Aufmerksamkeit
            \item Werbefunktion
        \end{itemize}
    \end{itemize}

    \subsubsection{erweitertes Produkt: Zusätzliche Dienstleistungen}
    \begin{itemize}
        \item Wichtig für Wettbewerb, weil diese Ebene viel Potential zur Differenzierung bietet.
    \end{itemize}

    \subsection{Produktlebenszyklus}
    Was impliziert Mahalos Position für die Marketingstrategien für Reisefanatiker 
    \begin{itemize}
        \item Produktdesign
        \item Pricing: Zahlbereite Kunden, nicht preissensibel: Skimming    
        \item Werbung: Geringes Budget weil Mahalo bereits ein bekanntes unternehmen ist gemäss Case-Study
    \end{itemize}

    \subsubsection{Preisuntergrenze}
    \begin{itemize}
        \item Um viel könnte man den Preis senken ohne Verluste zu realisieren?
        \begin{itemize}
            \item $Gewinn = U - K = p \cdot x - (k_{fix} + k_{var} \cdot x)$
            \item Bedingung kein Verlust: $U = K$
            \item Reduktion des Preises mit Unterschied (y): $p = (p_0 - y)$
            \item Wenn Absatz x reduziert kann es wieder in der Funktion eingesetzt werden
            \begin{itemize}
                \item Zum Beispiel wegen Konkurrenz nur noch 2/3 des alten Absatzes wird verkauft also $x_{alt} * \frac{2}{3}$ 
            \end{itemize}
        \end{itemize}
    \end{itemize}

    \subsubsection{Preisobergrenze oder Prohibitivpreis}
    Bei welchem Preis kann man nichts mehr verkaufen?
    \begin{itemize}
        \item Beim y-achsenschnitt der preis-mengen/absatz-Funktion
        \item p(x) mit x = 0
    \end{itemize}

    \subsubsection{Gewinnmaximierung}
    \begin{itemize}
        \item Am Punkt, wo die Gewinnfunktion 1. Ableitung = 0 und 2. Ableitung negativ 
    \end{itemize}

    \subsubsection{Preiselastizität}
    \[
        \epsilon = \frac{\delta x / x}{ \delta p / p} = \frac{dx *p}{dp * x} \\
        = \frac{}{}{\%Mengenänderung} {\%Preisänderung}
    \]
    \begin{itemize}
        \item $\epsilon > -1$
        \begin{itemize}
            \item Unelastisch (Preiserhöhung wirkt unterproportional auf Nachfrage => man kann preise erhöhen ohne Kunden zu verlieren)
        \end{itemize}
        \item $\epsilon < -1 $ (Preiserhöhung wirkt überproportional)
        \begin{itemize}
            \item Elastisch 
        \end{itemize}
    \end{itemize}

    
\end{document}  